\section{Introduction}
In-Context Learning (ICL) refers to the ability of a machine learning model to solve a task conditioned on examples of input-output pairs given in the prompt at inference time, without any updates to the model weights. So the "learning" (of weights $w$ for linear $f(x)=w^Tx$ for example) happens in the activations during the forward pass, while the weights of the model only encode the optimization algorithm.

An intuitive example motivating how tranformers can encode optimization procedures, such as gradient descent is presented by \cite{vonOswald2023Transformers}, restated here in a simpler and less rigourus way. 

Consider a least squares problem with training data $\{(x_i, y_i)\}_{i=1}^N$ where $y(x)=w^Tx$, with objective function
\begin{equation}
    L(w)=\frac{1}{2N}\sum_{i=1}^N ||w^Tx_i-y_i||^2.
\end{equation}

Consider the weight update in one step of gradient descent with step size $\eta$
\begin{equation}
    \Delta w
= -\eta \nabla_w L(w)
= -\eta \sum_{i=1}^{N} (w^T x_i - y_i)x_i.
\end{equation}

Also assume initial weights $w_0=0$, then

\begin{equation}
    w_1 = w_0 +\Delta w_0 
= \eta \sum_{i=1}^{N} y_i x_i.
\end{equation}

The prediction for the query based on the approximation of $w$ after a single step of gradient descent is
\begin{equation}
    \hat{y_q} = w_1^Tx_q = \eta \sum_{i=1}^{N} y_i x_i^Tx_q.
\end{equation}




